\section*{Chapter \ref{ch:ll:vector-instructions-and-data-parallel-code} --- Vector Instructions and Data-Parallel Code}

\begin{solution}{exo:ll:vector-instructions-and-data-parallel-code:1}
4 doubles, 8 integers of 32 bits, 32 bytes. $1\,500 = 46 \times 32 + 28$: 46 full blocks and one overlapping last block, 47
comparisons.
\end{solution}

\begin{solution}{exo:ll:vector-instructions-and-data-parallel-code:2}
$\mathtt{0x00041010}$ has bits 4, 12 and 18 set: delimiters at bytes 4, 12 and 18 of the block.
\end{solution}

\begin{solution}{exo:ll:vector-instructions-and-data-parallel-code:3}
After the first step the 16-bit lanes hold 20, 26, 9 and 25; after the second the 32-bit lanes hold 2\,026 and 925; the third
gives $2\,026 \times 10\,000 + 925 = 20\,260\,925$.
\end{solution}

\begin{solution}{exo:ll:vector-instructions-and-data-parallel-code:4}
Bytes 0--31 and 32--63, then the last load covers 38--69. After two full blocks the position is 64, so the last mask is shifted
right by $64 - 38 = 26$ bits: its remaining bits are bytes 64--69. A delimiter at byte 60 is found in the second block, before
the last load; the shift only discards bytes already examined.
\end{solution}

\begin{solution}{exo:ll:vector-instructions-and-data-parallel-code:5}
Possible aliasing: the output might overlap an input, so the loop is versioned with a run-time check; \texttt{\_\_restrict}
promises no overlap and removes it. No vector type for the record's fields: the fields are 80 bytes apart; store them as
columns. Control flow in the loop: the search may return early; write it with intrinsics, or search a whole block with a
mask as the build does.
\end{solution}

\begin{solution}{exo:ll:vector-instructions-and-data-parallel-code:6}
About \qty{52}{\nano\second} on the committed measurement, some 79 bytes per nanosecond. Once the buffer is beyond the
first-level cache, the rate at which the next levels deliver lines limits it, not the comparisons.
\end{solution}

\begin{solution}{exo:ll:vector-instructions-and-data-parallel-code:7}
Copy \texttt{positions\_any\_avx2} with 128-bit registers (\texttt{\_mm\_cmpeq\_epi8} and
\texttt{\_mm\_movemask\_epi8}) and 16-byte blocks (the last one at $n - 16$, shifted by $i - (n - 16)$), fall back to the scalar loop below 16 bytes, make the
dispatch return it at \texttt{Level::sse2}, and add a comparison with \texttt{positions\_any\_scalar} in the length loop of
the test.
\end{solution}

\begin{solution}{exo:ll:vector-instructions-and-data-parallel-code:8}
\texttt{-march=native} compiles for the build machine's processor, AVX-512 included, everywhere the compiler chooses to use
it; the continuous-integration runners and the new servers happened to have it too. The first older production host executes
an instruction it does not have and the process dies with an illegal-instruction signal. Compile for the oldest processor in
the fleet and dispatch wider code at run time.
\end{solution}

\begin{solution}{pb:ll:vector-instructions-and-data-parallel-code:1}
\begin{enumerate}
\item About \qty{0.22}{\nano\second} a byte for the scalar loop, \qty{0.024}{\nano\second} for SSE2 and \qty{0.013}{\nano\second}
for AVX2.
\item On the committed run, $a \approx \qty{0.3}{\nano\second}$ and $b \approx \qty{0.0125}{\nano\second}$ per byte.
\item $L^* = a/(b_{\text{scalar}} - b) \approx 2$ bytes for AVX2 and about 8 for SSE2.
\item The scalar loop also pays per call (the call, the exit branch), so its line lies above the one through the origin, and
the true crossing is earlier.
\item Scaling the measured split to 200 bytes: about \qty{120}{\nano\second} one byte at a time, \qty{14}{\nano\second} with AVX2.
\item $200\,000 \times \qty{120}{\nano\second} \approx \qty{24}{\milli\second}$ a second (2.4\% of a core) against
\qty{2.8}{\milli\second}.
\item $4 \times (42.6 - 6.6) \approx \qty{144}{\nano\second}$ per report.
\item The parsing: \qty{144}{\nano\second} against about \qty{106}{\nano\second} for the split.
\item There is nothing to search: fields are read at known offsets.
\item The dispatch calls the SSE2 version, which every x86-64 processor runs; nothing fails.
\item A load past the end of the message may cross into a page that is not mapped and fault, and it reads bytes that are not
the message's; the overlapping block reads only the message.
\item Twice the bytes per comparison, so roughly half the cost per byte on long buffers; nothing for short messages, where the
fixed cost dominates.
\item \textbf{Named result.} The AVX2 search costs about \qty{0.3}{\nano\second} per call plus \qty{0.0125}{\nano\second} per
byte against \qty{0.22}{\nano\second} per byte for the scalar loop: it repays its setup from about 2 bytes (SSE2 from about
8), so no message is too short for it.
\item Single fields can be (\texttt{54=1} is five bytes with its delimiter), but the scan runs over the whole message, never
over one field.
\item The same benchmark on the server's processor, with the real distribution of message lengths and with cold caches.
\item They define the correct answer, which the tests hold every vector version to, and they are the fallback.
\item Call each version directly in the tests, as the build does, and add a way to force the level (an environment variable
read by \texttt{detect}) so that the whole program can run on the fallback on a machine that has AVX2.
\item Checksums, copying, JSON and FIX parsing, searches over price levels, and bulk computations over positions.
\item Keep the fields the checks read (prices, quantities, limits) as columns.
\item When the same operation applies to many contiguous elements.
\end{enumerate}
\end{solution}

\begin{solution}{iq:ll:vector-instructions-and-data-parallel-code:1}
One instruction on several elements in a wide register. A loop adding two arrays of doubles element by element benefits;
walking a linked list, or a loop whose every iteration depends on the previous one's result, does not.
\iqlookfor{independent iterations on contiguous data.}
\end{solution}

\begin{solution}{iq:ll:vector-instructions-and-data-parallel-code:2}
Compare 16 or 32 bytes at once with the byte broadcast into a vector register, take the mask of equal bytes, and count its
trailing zeros; finish with an overlapping last block rather than reading past the end; choose the width at run time.
Check against a scalar loop on every length.
\iqlookfor{compare, mask, count, and the tail.}
\end{solution}

\begin{solution}{iq:ll:vector-instructions-and-data-parallel-code:3}
Possible aliasing, early exits, non-contiguous or strided data, calls that are not inlined, floating-point reductions whose
order the language fixes, and a cost model that judges it unprofitable. Ask the compiler: \texttt{-fopt-info-vec-missed} in
GCC gives the reason per loop; then read the assembly.
\iqlookfor{concrete blockers and the report flag.}
\end{solution}

\begin{solution}{iq:ll:vector-instructions-and-data-parallel-code:4}
Compile for the oldest processor, put the AVX2 code in functions compiled for AVX2 by attribute (or \omterm{def:ll:cpp-for-latency-iii-the-compiler:fmv}{function
multiversioning}), and choose at run time from the processor's reported features; test both paths.
\iqlookfor{run-time dispatch, not a global flag.}
\end{solution}

\begin{solution}{iq:ll:vector-instructions-and-data-parallel-code:5}
\omterm{def:ll:vector-instructions-and-data-parallel-code:soa}{Structure of arrays} for the revaluation: each Greek of consecutive positions is contiguous, so it fills vector registers and
fetches only the fields read. Records remain better for code that handles one position at a time; a hybrid of small blocks
serves both.
\iqlookfor{access pattern decides the layout.}
\end{solution}

\begin{solution}{iq:ll:vector-instructions-and-data-parallel-code:6}
Split at the decimal point, parse each part eight digits at a time inside a 64-bit register (subtract the character zero,
then combine pairs, quads and octets with three multiply-and-mask steps), and scale the fraction by the missing decimals.
Validate that every byte is a digit (a bit test on the whole word), the number of decimals, the range, and the sign.
\iqlookfor{the multiply-and-mask steps and the validation.}
\end{solution}
